\subsection{All completed toric pairs in the inherited cubic class-two quotient}

\begin{statement}{E.32}{Proved}{The complete convergence criterion, including the invisible inverse-limit obstruction}
Retain the arithmetic group $H\rtimes G_{\mathbb Q}$ of E.29.
For a set $I$ write
\[
 \mathcal R_I=\{(z_i)\in\textstyle\prod_{i\in I}\mathbb Z_3:
       \text{for every }n,\ z_i\in3^n\mathbb Z_3
       \text{ outside a finite subset of }I\}.
\]
Let $\mathcal P$ be the rational primes and
$\mathcal P^*=\mathcal P-\{3\}$. The Kummer isomorphism is
\[
 H^1_{\rm cts}(G_{\mathbb Q},\mathbb Z_3(1))
       =\mathcal R_{\mathcal P},\qquad
 (a_p)_p\longmapsto\sum_p a_p\kappa_p.                 \tag{E.32a}
\]
Choose a compatible system of roots for each $\kappa_p$;
the sum in (E.32a) converges uniformly as a cocycle.

Here is an explicit criterion for every pair
$a=(a_p),b=(b_p)\in\mathcal R_{\mathcal P}$.
For $q\ne3$ put $e_q=v_3(q-1)$ and choose a primitive
$3^{e_q}$th root $\eta_q\in\mathbb F_q^*$ when $e_q>0$.
For $p\ne q$ let $r_{q,p}\in\{0,\ldots,3^{e_q}-1\}$ be
specified by
\[
 \overline p^{(q-1)/3^{e_q}}=\eta_q^{r_{q,p}}.
\]
When $e_q=0$, put $r_{q,p}=0$. Define convergent sums and a
central correction by
\[
 A_q=\sum_{p\ne q}a_p r_{q,p},\qquad
 B_q=\sum_{p\ne q}b_p r_{q,p},\qquad
 c_q=-\frac{9(A_qb_q-B_qa_q)}{q-1}\in\mathbb Q_3.
                                                               \tag{E.32b}
\]
The fibre of $H^1_{\rm cts}(G_{\mathbb Q},H)$ over the pair
in (E.32a) has exactly three classes precisely when
\[
 c_q\in\mathbb Z_3\quad(q\ne3),\qquad
             (c_q)_{q\ne3}\in\mathcal R_{\mathcal P^*};       \tag{E.32c}
\]
in all other cases that fibre is empty. The criterion is
independent of the roots and generators used above.
For a nonempty fibre its three representatives are given by
(E.29d), starting with any one lift.

The continuous degree-two group, and the morphism used to
compute this obstruction, are
\[
 \begin{aligned}
 H^2_{\rm cts}(G_{\mathbb Q},\mathbb Z_3(2))
   &\xrightarrow[\sim]{\operatorname{loc}_{\rm cts}}
      \mathcal R_{\mathcal P^*}/d\mathcal R_{\mathcal P^*},\\
 9\Big(\sum_pa_p\kappa_p\Big)\cup
       \Big(\sum_pb_p\kappa_p\Big)
   &\longmapsto [\,9(A_qb_q-B_qa_q)\,]_q .
 \end{aligned}                                                 \tag{E.32d}
\]
Here $d(z)_q=(q-1)z_q$.
These are isomorphisms of $\mathbb Z_3$-modules; the denominator
is the actual image of $d$, with no closure added.
The map from this quotient to the product of local groups is
coordinatewise reduction modulo $q-1$. Its exact kernel is
\[
 \frac{\{z\in\mathcal R_{\mathcal P^*}:
                   z_q\in(q-1)\mathbb Z_3\text{ for all }q\}}
      {d\mathcal R_{\mathcal P^*}}.                            \tag{E.32e}
\]

This kernel contains an obstruction of an actual completed
toric pair. Choose recursively increasing primes $q_j$, $j\ge3$,
with
\[
 q_j\equiv1\pmod{3^j},\qquad
             3^{(q_j-1)/3}\not\equiv1\pmod{q_j},
\]
and write $r_j=r_{q_j,3}\in\mathbb Z_3^*$.
Take $a_3=1$, $a_p=0$ for $p\ne3$, and
\[
 b_{q_j}=\frac{q_j-1}{9r_j},\qquad b_p=0
              \text{ at the other primes}.                    \tag{E.32f}
\]
Then $b\in\mathcal R_{\mathcal P}$ and every local obstruction
and every finite-coefficient global obstruction vanishes.
There is nevertheless no continuous section with this completed
toric pair: its corrections are $c_{q_j}=-1$ for all $j$.
Every finite-support partial pair obtained from (E.32f) does
have three continuous section classes.

The classification concerns the inherited quotient $H$.
The map from the cubic's intrinsic quotient is exactly (E.29e);
the additional intrinsic obstruction E.31 remains present.
\emph{Manuscript: \S51.16, folios 733v--736v, scans 172--178,
(371)--(379), with its completed Kummer parameters retained.}
\end{statement}
\begin{proof}
We first calculate the continuous cohomology group and its
localization morphism, retaining the derived inverse limit.
Put
\[
 U=\operatorname{Spec}\mathbb Z[1/3],\qquad
 G_S=G_{\mathbb Q,\{3,\infty\}},\qquad
 T_n=\mathbb Z/3^n(2),\qquad T=\mathbb Z_3(2).
\]
Milne's comparison for locally constant sheaves
\cite[II.2.9]{E_MilneADT} identifies
$R\Gamma(U,T_n)$ with $R\Gamma(G_S,T_n)$.
We need the following actual vanishing:
\[
 H^2(G_S,T_n)=H^3(G_S,T_n)=0,
 \quad H^1(G_S,T_n)\text{ finite},\qquad n\ge1.                \tag{E.32g}
\]
For $n=1$, $T_1=\mathbb F_3$ with trivial action, since
$\chi_3^2\equiv1\pmod3$.
The Kronecker--Weber theorem
\cite[I.4.16]{E_MilneCFT} gives
$H^1(G_S,\mathbb F_3)=\mathbb F_3$.
More explicitly, an abelian extension unramified away from
three is contained in a cyclotomic field, and its quotient
kills the inertia factors at all other primes. Its cubic
quotient therefore factors through
$\operatorname{Gal}(\mathbb Q(\zeta_{3^m})/\mathbb Q)$,
which has a unique quotient of order three. The nontrivial
quotient already occurs in $\mathbb Q(\zeta_9)^+$.
Both $H^0(G_S,\mathbb F_3)$ and $H^1(G_S,\mathbb F_3)$
have order three. The global Euler characteristic
\cite[I.5.1]{E_MilneADT} equals one here: its real factor is
$|H^0(G_{\mathbb R},\mathbb F_3)|/3=1$.
Consequently $H^2(G_S,\mathbb F_3)=0$.
Filter $T_n$ by its powers of three. Each successive quotient
is the same trivial $\mathbb F_3$, so the long exact
cohomology sequence proves $H^2(G_S,T_n)=0$ for all $n$.
Milne I.4.10(c) identifies every degree at least three with
real-place cohomology, which is zero for a module of odd order.
Finiteness is I.4.15. This proves (E.32g).

For a finite set $J\subset\mathcal P^*$, apply the
localization triangle to
$U_J=U-\{q:q\in J\}$, using
\cite[II.0.1(f), II.1.7(b)]{E_MilneADT}.
Here the local purity identification can also be checked
directly: over the strict Henselization at $q$, a locally
constant $T_n$ is constant; the fraction field has tame
$3$-inertia $\mathbb Z_3(1)$, while its wild inertia is
pro-$q$. Its cohomology is $T_n$ in degree zero,
$T_n(-1)$ in degree one, and zero above degree one.
The local support complex is therefore
$T_n(-1)[-2]$, with its residue-field action.
Thus there is the compatible distinguished triangle
\[
 R\Gamma(U,T_n)\longrightarrow R\Gamma(U_J,T_n)
 \longrightarrow\bigoplus_{q\in J}
       R\Gamma(\mathbb F_q,\mathbb Z/3^n(1))[-1]
 \longrightarrow R\Gamma(U,T_n)[1].                         \tag{E.32h}
\]
First take the filtered direct limit over finite $J$.
Filtered direct limits of abelian groups are exact. On the
generic-field term this gives $R\Gamma(G_{\mathbb Q},T_n)$:
every continuous cochain with finite coefficients factors
through a finite quotient of $G_{\mathbb Q}$, whose number
field is ramified at only finitely many primes.
In particular the third term becomes a direct sum at each
fixed coefficient level $n$; no infinite product is substituted.

Next take $R\varprojlim_n$ of these triangles. The comparison
with continuous cochains is Jannsen's theorem
\cite[(2.2), together with (1.6)]{E_JannsenContinuous}.
For the first term (E.32g), the finiteness of the degree-one
groups kills their first derived inverse limit. It follows that
\[
 H^2(R\varprojlim_nR\Gamma(U,T_n))
 =H^3(R\varprojlim_nR\Gamma(U,T_n))=0.                       \tag{E.32i}
\]
For the support term, the procyclic residue Galois group has
the compatible two-term cochain model
\[
 R\Gamma(\mathbb F_q,\mathbb Z/3^n(1))
  \simeq[\,\mathbb Z/3^n\xrightarrow{q-1}\mathbb Z/3^n\,]
                 \quad\text{in degrees }0,1.                \tag{E.32j}
\]
Choose arithmetic Frobenius for this model. After taking the
direct sum over $q\ne3$, the transition in each of its two
terms is the surjective reduction
$\bigoplus_q\mathbb Z/3^{n+1}\to\bigoplus_q\mathbb Z/3^n$.
The inverse systems of terms are therefore acyclic for
$\varprojlim$. Their derived inverse limit is the literal
complex
\[
 [\,\mathcal R_{\mathcal P^*}
         \xrightarrow{d}\mathcal R_{\mathcal P^*}\,].          \tag{E.32k}
\]
Indeed a compatible sequence of finitely supported residues
is exactly a family of $3$-adic integers tending to zero.
The degree-two segment of the resulting triangle, using
(E.32i), proves the first isomorphism in (E.32d).
This calculation takes the inverse limit of complexes. It
does not replace it by the inverse limit of their cohomology.

We next compute the image of the actual obstruction under
this morphism. The Kummer sequence and
$\mathbb Q^*=\{\pm1\}\times\bigoplus_p p^{\mathbb Z}$
give (E.32a); the sign factor has trivial pro-$3$ completion.
Multiplying a fixed cocycle $\kappa_p$ by $a_p$ and summing
is legitimate on continuous cochains: for each $n$, all but
finitely many summands vanish modulo $3^n$, uniformly on
$G_{\mathbb Q}$. Write these sums as $\alpha,\beta$.
Their obstruction is the continuous cocycle
\[
 \omega(g,h)=\frac92\{
    \alpha(g)\chi_3(g)\beta(h)
       -\beta(g)\chi_3(g)\alpha(h)\},                        \tag{E.32l}
\]
representing $9\alpha\cup\beta$ as in E.29.

Fix $q\ne3$ and a decomposition group. Choose tame inertia
$\tau_q$ and arithmetic Frobenius $\phi_q$ such that
\[
 \phi_q\tau_q\phi_q^{-1}=\tau_q^q,
 \qquad \kappa_q(\tau_q)=1,
 \qquad \kappa_q(\phi_q)=0.                                 \tag{E.32m}
\]
These choices use compatible roots of a uniformizer, as in
E.29. Put
$\alpha(\tau_q)=a_q$, $\beta(\tau_q)=b_q$ and
$\alpha(\phi_q)=A'_q$, $\beta(\phi_q)=B'_q$.
On the full inertia group the two cocycles are proportional
to the same tame character; wild inertia maps to zero.
Consequently (E.32l) is identically zero on two inertia
arguments, giving its canonical central-zero inertia splitting.
In the pullback of $H\rtimes G_{\mathbb Q}$, let the central
coordinates of the two lifts be $w$ at $\tau_q$ and $z$ at
$\phi_q$. Direct multiplication in (E.29a) gives the defect
between the two sides of (E.32m) as
\[
 q\{(q-1)w+9(A'_qb_q-B'_qa_q)\}.                            \tag{E.32n}
\]
It is independent of $z$. In particular the central-zero
inertia splitting represents the residue cochain with value
$9(A'_qb_q-B'_qa_q)$ at Frobenius. The division by $q$ in
(E.32n) is the identification
$\operatorname{Hom}(\mathbb Z_3(1),T)=T(-1)$:
the conjugation defect at $\tau_q^q$ is pulled back to its
value at $\tau_q$. Changing the inertia splitting by $w$
adds precisely $(q-1)w$, the differential in (E.32j).
This is the chain-level localization morphism in (E.32h):
it records the failure of an inertia splitting to be invariant
under residue-field Frobenius.

At each finite coefficient level these inertia splittings
agree with an unramified splitting outside finitely many
primes, since $\alpha,\beta$ factor through a finite quotient
modulo $3^n$. They therefore form a compatible representative
in the direct-sum complexes used in (E.32h)--(E.32k).
The resulting family tends to zero, since it contains a
factor $a_q$ or $b_q$ in each summand.
This establishes the localization representative itself,
including its class in the quotient (E.32k).

For a unit $p$ at $q$, finite-field Kummer theory gives
\[
 \kappa_p(\phi_q)\equiv r_{q,p}\pmod{q-1}.
\]
Choose the place over $q$ to send the selected global
$3^{e_q}$th root of unity to $\eta_q$.
The formula follows by applying Frobenius to a
$3^{e_q}$th root of $p$: its quotient is
$\overline p^{(q-1)/3^{e_q}}$.
Write
$\kappa_p(\phi_q)-r_{q,p}=(q-1)t_{q,p}$ with
$t_{q,p}\in\mathbb Z_3$.
The sums of $a_pt_{q,p}$ and $b_pt_{q,p}$ converge and are
uniformly bounded in $\mathbb Z_3$.
Hence replacing $A'_q,B'_q$ by $A_q,B_q$ changes the
representative in (E.32d) by $d$ of a family tending to zero:
the change divided by $q-1$ has a factor $a_q$ or $b_q$.
This proves the second formula in (E.32d).
The same argument proves independence of root choices.
Changing a residue generator scales the corresponding
coordinate by a $3$-adic unit and adds a term of this form,
so it also preserves (E.32c).

For completeness, changing global cocycles by
$\alpha\mapsto\alpha+du$, $\beta\mapsto\beta+dv$, with
$du(g)=(\chi_3(g)-1)u$, changes the residue representative by
\[
 (q-1)\,9(ub_q-va_q).
\]
The family in the second factor lies in
$\mathcal R_{\mathcal P^*}$.
The corresponding central inertia correction changes by
$-9(ub_q-va_q)$, exactly as the change
$z_g\mapsto h^{-1}z_g g(h)$ with $h=(u,v,0)$ gives.
Also replacing $\phi_q$ by
$\tau_q^r\phi_q$ leaves $A'_qb_q-B'_qa_q$ unchanged.
These formulas specify both the coboundary sign and the
actual coefficient comparisons.

The obstruction vanishes precisely when its representative
lies in $d\mathcal R_{\mathcal P^*}$. Since $q-1\ne0$ in
$\mathbb Z_3$, its unique possible central correction is
(E.32b). This proves (E.32c) and the empty-fibre assertion.
The central group has
$H^1_{\rm cts}(G_{\mathbb Q},T)=\mathbb Z/3$ by E.29,
and $H^0(G_{\mathbb Q},\mathbb Z_3(1)^2)=0$.
The exact nonabelian central-extension sequence therefore
makes each nonempty fibre a torsor under this group.
Its three explicit twists are (E.29d).

We finally verify the inverse-limit example. The field
$\mathbb Q(\sqrt[3]{3})$ cannot be contained in a cyclotomic
field: a subfield of an abelian Galois extension is Galois,
whereas this irreducible real cubic does not contain the two
nonreal roots of $X^3-3$.
Thus
$\mathbb Q(\zeta_{3^j},\sqrt[3]{3})/
 \mathbb Q(\zeta_{3^j})$
has degree three. A nonidentity element of this subgroup has
a conjugacy class restricting trivially to the cyclotomic
field. Chebotarev
\cite[V.3.23]{E_MilneCFT} provides infinitely many rational
primes with that class, giving the two requirements on $q_j$.
Choose each $q_j$ to be the least such prime larger than the
previous choice. Noncubicity says $r_j$ is a $3$-adic unit,
so $v_3(b_{q_j})=e_{q_j}-2\ge j-2$.
This proves that (E.32f) is an actual completed Kummer pair.
Its representative in (E.32d) is
\[
 z_{q_j}=q_j-1,\qquad z_q=0\text{ otherwise}.                \tag{E.32o}
\]
It is a restricted family, but its only preimage under
coordinatewise multiplication by $q-1$ is the family with
value one at every $q_j$, which is not restricted.
Thus its class is nonzero in (E.32d), despite lying in the
kernel (E.32e). At every finite level its class is zero in
the cohomology of the direct-sum complex: only finitely
many entries of (E.32o) survive modulo $3^n$, and they
are obtained by applying $d$ to the finitely supported
family with value one at those entries.
The finite-level global obstruction is zero by
(E.32g)--(E.32h). At $q\ne3$ the local obstruction is zero
by coordinatewise divisibility. At three, finite local
duality \cite[I.2.3]{E_MilneADT} gives
$H^2(\mathbb Q_3,T_n)=H^0(\mathbb Q_3,\mathbb Z/3^n(-1))^*=0$:
the local cyclotomic character contains $-1$, whose action
has no fixed vector on this odd-order module.
Finite local degree-one groups then give
$H^2_{\rm cts}(\mathbb Q_3,T)=0$ as well.
Real cohomology with these coefficients is zero.
Every finite-support partial pair in (E.32f) has finitely
many corrections equal to $-1$, so it satisfies (E.32c).
All assertions follow.
\end{proof}
